\documentclass[11pt]{article}
\usepackage{amsmath,amssymb,amsthm,geometry}
\geometry{margin=.92in}\setlength{\emergencystretch}{1.5em}
\newtheorem{theorem}{Theorem}
\title{Retention of pre-weak probe depth through weak learning}
\author{AN04, E5 step 5}\date{Version 1}
\begin{document}\maketitle
\paragraph{Question and conventions.}
Normalize time as in P. Use the aligned balanced isotropic critical
family $S_0=cq^{1/\omega}$, $c>0$ fixed,
$\lambda=\rho^2$, $\omega=(1-\lambda)/2$. For the benchmark choose
$(\rho,\zeta)$ on WIT's fixed learned rectangle and put
$e_*=\min_{s\in J}\mathcal E_{\rm coh}(s)$. It is finite and negative:
on $J$, $A=\zeta-X>0$ by WIT uniformtheorem and $Z=X-Y>0$
by AN01 barrier sign, so the coherent energy equals $-AZ-A^2/2$.
Write $\xi=(\sin(q\zeta),-\cos(q\zeta))$, $r_t(z)=z-f_t(z)$ and
$e_q(t)=(|r_t(\xi)|^2-1)/(2q^2)$.
The strong half-mass clock $t_h$ is that of HBD; the weak-response
hit $t_{1/2}$ is that of BRW. For a concrete measured target, let
$\mathcal L_q$ consist of the times in $t_{1/2}+J$ satisfying the
two response bounds $\mathsf m_1\ge1/2,\mathsf m_2\ge1/4$ and
the two cluster-loss ratios at most $169/196$. If there is no hit,
take the set empty. These are evaluation criteria; RTA's attainment
theorem applies to the ten-power family. Other learned windows or a
fixed weak-response test can be used in the transport identity below.
A clock hit alone would not prove the cluster-loss guards.

\section{How much of the depth does transfer?}
\begin{theorem}[Trained, before weak learning]\label{an04drrv1:interval}
There are delays $s(q)\to\infty$ and lengths $w(q)\to\infty$ such
that, on $I_q=t_h+[s(q),s(q)+w(q)]$,
\[
 \sup_{I_q}e_q\to-\infty,\qquad
 \sup_{I_q}|f_t(\xi)|\to0,\qquad
 \sup_{[t_h,\sup I_q]}|\mathsf m_2|\to0 .
\]
No explicit growth rate for the delays or lengths is asserted.
\end{theorem}
\begin{proof}
COD main and trained give the common limiting upper bound
$C-k e^{\lambda s/2}/(1+s)$ at fixed $s$. In fact on any fixed compact
interval of delays the same limsup bound applies at moving times
$s_q\to s$ with the bound evaluated at $s$. To justify this step,
BRW extracts right interior label paths uniformly in time on each
such interval. COD's two scaled probe integrands on these compacts
are continuous positive-part functions, converging uniformly on
their bounded charts. Its physical bound gives an $O_I(\delta)$
contribution from axis strips of label length $\delta$ to both
$f_1/q$ and $f_2/q^2$, uniformly on that time interval. The Q2
compact is probe-inactive by COD's two positive component bounds
through the endpoint; the compact Q3 mass remains $O_I(S_0)=o(q^2)$
throughout. Thus COD's passage to the probe works also at $s_q$.
In particular both scaled output components are bounded on a fixed
interval and $|f_t(\xi)|=O_I(q)$ there.
For integer $j$ choose increasing $s_j\ge j$ with the common energy
bound below $-2j$ for every $s\ge s_j$. Apply the compact assertion
on $[s_j,s_j+j]$ and BRW delay on $[0,s_j+j]$. Choose $q_j\downarrow0$
so that for $q\le q_j$ the supremum of $e_q$ on the former interval
is below $-j$, that of $|f_t(\xi)|$ below $1/j$, and that of the
weak response in absolute value on the latter interval below $1/j$.
On $q_{j+1}<q\le q_j$ choose $s(q)=s_j$, $w(q)=j$.
The selected intervals end before any fixed positive weak-response
threshold for small $q$. The assertion supplies no bounds at that hit.
\end{proof}

\section{Exact retention test with weak feedback}
Take $t_0=t_h+s_K$ to be a COD witness with $e_0=e_q(t_0)<e_*$,
fixing $s_K$ before $q$, or take any witness from the diagonal
interval above. Set $r_0=r_{t_0}(\xi)$, $R_0=|r_0|$ and
$\hat r=r_0/R_0$. For the witnesses under consideration $R_0\to1$.
Define the time-dependent probe row of the full tangent kernel by
\[
 k_t(\xi,D)=\int\left[
 (a^T\xi)_+(a^TD)_+ I_2+
 (\xi^TD){\bf1}_{a^T\xi>0}{\bf1}_{a^TD>0}\,bb^T
 \right]\,d\mu_t .
\]
Here $a,b$ are the input/output unit directions of a current
balanced label, $d\mu_t=m\,d\alpha/(2\pi)$ and the mixture expectation
is over the two training Gaussians. For finite $T\ge t_0$ put
\begin{align}
 G(T)&=\int_{t_0}^{T}\mathbb E[k_t(\xi,D)r_t(D)]\,dt,\nonumber\\
 B_q(T)&=q^{-2}\hat r^T G(T),\qquad
 A_q(T)=q^{-1}(\hat r^\perp)^T G(T). \label{an04drrv1:projections}
\end{align}
The letters in this display are signed changes in probe output, not
radial masses or the strong exposure. This fixes a direction at the
deep witness; it does not keep the evolving network or its residual
fixed. Differentiating the two physical factors of each ReLU
feature by P exactflow gives $\dot f_t(\xi)=\mathbb E[k_t r_t]$
almost everywhere, as in BRW's kernel derivation. The positive-part
chain rule permits gate crossings. Thus $G=f_T(\xi)-f_{t_0}(\xi)$.
Polarization at $t_0$ gives the exact identity
\begin{equation}\label{an04drrv1:test}
 e_q(T)-e_0=\tfrac12 A_q(T)^2+\tfrac12q^2 B_q(T)^2-R_0 B_q(T).
\end{equation}
It follows that the missing strict inequality, in independently
specified field integrals, is
\begin{equation}\label{an04drrv1:missing}
 \inf_{\substack{T\in\mathcal L_q\\T\ge t_0}}
 \left\{\tfrac12 A_q(T)^2+\tfrac12q^2 B_q(T)^2-R_0 B_q(T)\right\}
       < e_*-e_0 .
\end{equation}
The infimum of the empty set is $+\infty$. For small $q$ any reached
learned window lies after the fixed witness by BRW delay; for the
diagonal construction the $1/j$ response bound and the fixed width
of $J$ also exclude earlier learned times for small $q$.
Strict inequality is equivalent to beating the benchmark somewhere
on the chosen learned window. A fixed positive asymptotic gap
requires a fixed strict margin. In particular a negative longitudinal
projection consumes the initial gain, and a transverse projection
of either sign incurs a square cost. Step 4 bounds the right side
from below but does not sign these two projections through a positive
weak-response hit. The following two checks locate actual failures
of attempts to deduce that bound.

\paragraph{First attempt: global movement budget (trained flow).}
Write $U,W$ for the physical vectors, $S=\int m\,dl$,
$dl=d\alpha/(2\pi)$ and
$g^2=\int(|\dot U|^2+|\dot W|^2)\,dl$. Direct differentiation of
$f_t(z)=\int W(U^Tz)_+dl$ and Cauchy--Schwarz give
$|\dot f_t(z)|\le |z|\sqrt{2S}\,g$. The same absolute budget
$P_T=\int_{t_0}^T\sqrt{2S}\,g\,dt$ thus bounds $|G|$.
A sufficient sign-blind estimate would be
\begin{equation}\label{an04drrv1:falsebudget}
 P_T<\sqrt{1+2q^2e_*}-R_0
 =\frac{2q^2(e_*-e_0)}{\sqrt{1+2q^2e_*}+R_0}.
\end{equation}
The right side is $O_{s_K}(q^2)$ for a fixed witness and is $o(1)$
even on the deep diagonal by Theorem~\ref{an04drrv1:interval}.
But from the exact weak response definition under $P_2$,
\[
 |\mathsf m_2(T)-\mathsf m_2(t_0)|
 \le\frac{\mathbb E_{P_2}|D_2||D|}{\rho^2+q^2}P_T
 \le\sqrt{\frac{\rho^2+2q^2}{\rho^2+q^2}}P_T .
\]
At every learned time $P_T\ge(1/4-o(1))/\sqrt2$.
Consequently \eqref{an04drrv1:falsebudget} is false on any
reached learned interval for small $q$, even if dissipation were
estimated exactly. This lower bound is on the common majorant,
not on the norm or the sign of the actual probe change.

\paragraph{Second attempt: forward invariance of depth (leading system).}
There is an explicit self-consistent flow test for the proposed rule
that a sufficiently deep pre-weak state must stay below $e_*$ until
weak learning. Use P \texttt{pa:6d} with one atom in each family,
strong mass one, actual weak mass $n'=\lambda n(1-n)$ and coordinates
$x,y,u,v$. Given any $K_0>0$, start at
$x(0)=0,\ y(0)=-R,\ u(0)=0$, where $R>0$ is chosen so that
$-\zeta^2/2-R\zeta<\min(-K_0,e_*)$.
Choose $0<n_0<1/4$ sufficiently small that the first time
$t_{1/4}$ of $n=1/4$ is after time one and the matched window
$t_{1/2}^{n}+J$ also starts after one. Here
$t_{1/2}^{n}=\lambda^{-1}\log[(1-n_0)/n_0]$.
Fix a finite $T$ beyond that window and $t_{1/4}$, and choose $v(0)$
last. All data are prescribed in the leading system.

The doubled velocities are
\begin{align*}
 2x'&=\lambda\Phi(\rho x)[nv-x+(1-n)y],&
 2y'&=-y-nK(u)+\rho(1-n)K(\rho x),\\
 2u'&=(1-n)(\phi(u)-\lambda u)-(nu+y)\Phi(u),&
 2v'&=\rho K(\rho x)-\lambda v .
\end{align*}
They have at most linear growth and are smooth. To control the
unfavorable core output, set $H=y_-+u_+$. Wherever $y\le0$,
$2(y_-)'\le K(u)\le\phi(0)+u_+$, with upper right derivative
at zero. Wherever $u\ge0$, $2(u_+)'\le\phi(0)+y_-$, also at
zero. On the other sides the corresponding positive-part velocity
vanishes. Thus $D^+H\le\phi(0)+H/2$ and
$H\le H_T=(R+2\phi(0))e^{T/2}$ through $T$, independently of $v(0)$.
Choose $d=\zeta+1$ and a positive $v(0)$ such that
\[
 n_0v(0)e^{-\lambda T/2}>d+H_T+8d/\lambda .
\]
Since $K\ge0$, $v(t)\ge v(0)e^{-\lambda t/2}$. Throughout
$0\le x\le d$ one has $\Phi(\rho x)\ge1/2$, hence $x'>2d$.
At the initial lower face $x=0$ the velocity is upward, and the
upper endpoint is a crossing face with the same strict sign.
The input reaches $d$ before time $1/2$; the face $x=d$ forbids
downward return through $T$. The leading probe error is exactly zero
after clearing on this interval since both atoms are inactive.
In particular the entire matched window has error $0>e_*$; weak
mass passes through $1/4$ and $1/2$ under its own equations.
The core state and both weak angles evolve throughout.

The same example supplies a strict outward contact for the
depth sublevel if desired. On $0\le x<\zeta$, write
$A=\zeta-x$, $Z=x-y$. On $\mathcal E=e_*$,
$Z=-e_*/A-A/2$. For an immediate boundary test choose any
$0<A<\min(\zeta,\sqrt{-2e_*})$, set $x=\zeta-A,\ y=x-Z$,
and take $u=0$, $0<n<1/4$. Then $Z>0$, and the full derivative
is $Zx'+Ay'$. Its coefficient of $v$ is
$\lambda Z\Phi(\rho x)n/2>0$. All other terms at this state are
independent of $v$. Thus a large finite $v$ gives a strictly outward
velocity on the depth face. This refutes the unrestricted sublevel
barrier. The example is neither a critical initialized trajectory
nor a counterexample to E5 for that family; in particular its
weak output tilt was freely prescribed. Applying this test to the
trained entrance would require bounding the actual weak feedback.

\paragraph{Finite-noise check for the prescribed entrance.}
The same loss of depth is possible in the original Gaussian
equations with the prescribed balanced two-atom initial state.
Fix all its constants, including $v(0)$ and $T$, before $q$.
Lift the strong vectors to angles $qx,qy$ and the weak vectors to
$\pi/2-qu,\pi/2+qv$, with radial masses $m=1,n=n_0$ initially.
The leading trajectory is bounded on a slightly extended finite interval.
Here is the compact field check.

Write $D=(1+qZ_1,qZ_2)$ on $P_1$ and $(qZ_1,\rho+qZ_2)$ on
$P_2$. On bounded mass and angular boxes, the components
of $r=D-f(D)$ are
\[
\begin{array}{c|cc}
 &r_1&r_2\\ \hline
 P_1&(1-m)D_1+O(q^2)&
 q[Z_2-my-n(Z_2+u)_+]+O(q^2)\\
 P_2&q[Z_1-m(Z_1+\rho x)_++\rho nv]+O(q^2)&
 (1-n)D_2+O(q^2).
\end{array}
\]
The errors hold in each fixed finite Gaussian $L^p$ norm, using
Lipschitz continuity of positive part and the own-cluster Gaussian
tail. In the exact angular projections $b^Tr(a^\perp\cdot D)$
times the receiver gate and $(b^\perp)^Tr(a^TD)_+$, divide by
$q$ and negate the weak input row for the $u$ convention.
The resulting doubled angular field has the following cluster limits,
where $P=\Phi(\rho x)$ and $Q=\Phi(u)$:
\[
\begin{array}{c|cc}
 &P_1&P_2\\\hline
 2x'&-(1-m)x&\rho(1-m)\phi(\rho x)+
          \lambda[nv-mx+(1-n)y]P\\
 2y'&-y-nK(u)&\rho(1-n)K(\rho x)\\
 2u'&(1-n)\phi(u)-[nu+my+(1-m)v]Q&-\lambda(1-n)u\\
 2v'&-(1-m)K(u)&\rho mK(\rho x)-\lambda v
\end{array}
\]
Each row sums to the doubled velocity. For example on $P_2$ the
strong cross gate tends to $\mathbf1_{Z_1+\rho x>0}$, and the
divided projected residual in its input row tends to
$Z_1-m(Z_1+\rho x)_++\rho[nv+(1-n)y]$.
Use $\mathbb E[Z{\bf1}_{Z+w>0}]=\phi(w)$; the $P_1$ weak cross
gate uses $Z_2+u$ in the same way. The log radial rates tend to
$1-m,\lambda(1-n)$ by the undivided projections. These convergences
are uniform on every fixed compact box: for any convergent sequence
of its parameters the cross gates converge almost surely outside
one Gaussian level set, with uniformly integrable divided products
as given by the residual expansions. The own gates have exponentially
small complements. Compactness gives uniformity. The leading field
is locally Lipschitz; Gronwall up to a first exit from a box strictly
enclosing its path excludes that exit and proves uniform finite-time
convergence for this prepared initial state.

Consequently the original weak response converges uniformly to $n$,
and the strong response to $1$, by the same residual table. The first
weak half-response hit converges to $t_{1/2}^{n}$ by strict logistic
crossing and uniform convergence. The strong loss tends uniformly
to zero, and the weak loss to $\rho^2(1-n)^2/2$. On the matched
window WIT gives $n\ge1/4+\sigma$ for some $\sigma>0$.
Thus these losses clear the numerical cutoffs against the critical
family's isotropic baseline losses (which tend to $1/2,\rho^2/2$).
The test itself has a different initialization. Choose $T$ with a
strict time margin as above. The strong gate stays inactive on the
full measured window since the limiting $x\ge\zeta+1$ there. The
weak gate is inactive by its bounded chart. The exact probe prediction
is then zero and $e_q=0$ throughout that window. Initially its
scaled energy tends to $-\zeta^2/2-R\zeta<-K_0$. This verifies a
failure of the depth-only implication even for the exact finite-noise
vector field, while leaving the critical trained family's sign undecided.

\paragraph{Disposition and handoff.}
The new positive result transfers deep pre-weak error over a growing
but unspecified interval. The learned-depth comparison is open.
Two separate quantities remain unpaid for the critical initialization:
existence and alignment of a learned window, and the projected
probe motion \eqref{an04drrv1:missing} on at least one reached
time there. The kernel expectations in \eqref{an04drrv1:projections}
give a target for a weak-passage estimate that avoids paying for
large training motion invisible to this probe. The example identifies
$nv$ in the strong input equation as one concrete term that a
gate-retention argument would have to control. It does not establish
its size on a trained learned window.

\paragraph{Source ranges.}\raggedright
COD \texttt{an04codv1:main, trained} and BRW \texttt{an04brwv1:delay}
concern the specified trained critical family only at the respective
fixed resident times or pre-weak delay. Compact-window extraction
uses their physical endpoint estimates under the same initialization.
P \texttt{pa:exactflow, learning, rawvelocity} use the finite-noise
Gaussian population, finite moments and the a.e.\ probe chain rule.
P \texttt{pa:6d} concerns prescribed atomic leading states;
WIT \texttt{inc:AN06:v2:uniformtheorem} and
AN01 \texttt{an01pbv1:sign} concern the selected branch and common
learned window for the benchmark only. No small-noise critical
weak-seed law, trained rate transfer through weak learning, or
numerical premise is used.
\end{document}
